\diarytitle{0034}{定数係数2階線形非同次微分方程式の初期値問題（特性方程式が重根で二重に共鳴する場合）}

未定係数法を用いて特殊解を求める。定数係数2階線形非同次方程式 $y'' + a y' + b y = f(x)$ の一般解は、対応する同次方程式の一般解 $y_h$ と、非同次方程式の特殊解 $y_p$ の和
\[
y = y_h + y_p
\]
で与えられる。特殊解は $f(x)$ の形に応じて多項式・指数関数・三角関数などの仮定形を置き、方程式に代入して未定係数を決定する。ただし仮定した形が同次解と重なる場合（共鳴）は、仮定形に $x$ の適当なべきを掛けて同次解と一次独立にする必要がある。

\medskip
\noindent
\textbf{同次解:} 特性方程式 $r^{2}-2r+1 = (r-1)^{2} = 0$ より重根 $r=1$。よって
\[
y_h = (C_{1}+C_{2}x)e^{x}
\]

\noindent
\textbf{特殊解（重根との共鳴）:} $f(x)=e^{x}$ は $e^{x}$ 型であるが、$e^{x}$ は同次解に（重複度2で）含まれるため、単純な $Ae^x$ でも $Axe^x$ でも共鳴してしまう。そこで
\[
y_p = Ax^{2}e^{x}
\]
とおく。
\[
y_p' = A e^{x}(x^{2}+2x), \qquad y_p'' = A e^{x}(x^{2}+4x+2)
\]
したがって
\[
y_p'' - 2y_p' + y_p = Ae^{x}\bigl[(x^{2}+4x+2) - 2(x^{2}+2x) + x^{2}\bigr] = Ae^{x}\cdot 2 = 2Ae^{x}
\]
これが $e^{x}$ に等しいので $2A=1$、すなわち $A = \dfrac{1}{2}$。よって $y_p = \dfrac{1}{2}x^{2}e^{x}$。

\noindent
\textbf{一般解と初期条件:}
\[
y = (C_{1}+C_{2}x)e^{x} + \frac{1}{2}x^{2}e^{x} = e^{x}\left(C_{1} + C_{2}x + \frac{1}{2}x^{2}\right)
\]
\[
y' = e^{x}\left(C_{1} + C_{2}x + \frac{1}{2}x^{2}\right) + e^{x}(C_{2}+x)
\]
$y(0) = C_{1} = 1$。$y'(0) = C_{1}+C_{2} = 0$ より $C_{2} = -1$。

\noindent
したがって解は
\[
\boxed{\; y = \left(1 - x + \frac{1}{2}x^{2}\right)e^{x} \;}
\]

（検算: $y' = e^x\cdot\dfrac{1}{2}x^2$、$y'' = e^x\left(\dfrac{1}{2}x^2+x\right)$ より
$y''-2y'+y = e^x\left[\left(\dfrac{1}{2}x^2+x\right) - 2\cdot\dfrac{1}{2}x^2 + \left(1-x+\dfrac{1}{2}x^2\right)\right] = e^x\cdot 1 = e^x$ OK。
また $y(0)=1$, $y'(0)=\dfrac{1}{2}\cdot 0^2=0$ OK）
